\chapter{Chapter09 — Coalitional Influence Asymptotics}
%%% generated by the blueprint pipeline from FABL.Chapter09.CoalitionalInfluenceAsymptotics
\section{Overview}

Declarations of the Lean module \texttt{FABL.Chapter09.CoalitionalInfluenceAsymptotics}.

\section{Declarations}

\begin{theorem}[Influence Bounded by Complementary Output Probability]
\lean{FABL.coalitionalInfluenceToward_le_complement_probability}
\label{FABL.coalitionalInfluenceToward_le_complement_probability}
\leanok
\uses{FABL.BooleanFunction, FABL.Sign, FABL.coalitionalInfluenceToward, FABL.uniformProbability, FABL.uniformProbability_le_one}
\difficulty{3}
Let $n$ be a nonnegative integer, let $f:\{-1,1\}^n\to\{-1,1\}$, let $J$ be a set of
coordinates, and let $b\in\{-1,1\}$. The influence of $J$ toward $b$, defined as the
fraction of inputs from which $J$ can make $f$ attain $b$ minus the fraction of inputs
where $f$ already equals $b$, is at most one minus the fraction of inputs where $f$
equals $b$.
\end{theorem}

\begin{definition}[Tribes block coordinate embedding]
\lean{FABL.tribesBlockEmbedding}
\label{FABL.tribesBlockEmbedding}
\leanok
For natural numbers $w$ and $s$ and a block index $i \text{ in } \text{Fin}(s)$, the map
from $\text{Fin}(w)$ to $\text{Fin}(sw)$ sends each within-block coordinate $o$ to the
coordinate corresponding to the pair $(i,o)$ under the standard product-coordinate
equivalence. This map is injective.
\end{definition}

\begin{definition}[Tribes block coordinates]
\lean{FABL.tribesBlockCoordinates}
\label{FABL.tribesBlockCoordinates}
\leanok
\uses{FABL.tribesBlockEmbedding}
For natural numbers $w$ and $s$ and a block index $i \text{ in } \text{Fin}(s)$, the
block-coordinate set is the image of all within-block coordinates under the embedding
into $\text{Fin}(sw)$.
\end{definition}

\begin{theorem}[Cardinality of a block of coordinates]
\lean{FABL.card_tribesBlockCoordinates}
\label{FABL.card_tribesBlockCoordinates}
\leanok
\uses{FABL.tribesBlockCoordinates}
\difficulty{1}
Let $w$ and $s$ be natural numbers, and choose one of the $s$ blocks in a set of $sw$
positions arranged in $s$ blocks of $w$ positions each. The set of positions in that
block has cardinality $w$.
\end{theorem}

\begin{theorem}[Membership in a tribe block of coordinates]
\lean{FABL.mem_tribesBlockCoordinates}
\label{FABL.mem_tribesBlockCoordinates}
\leanok
\uses{FABL.tribesBlockCoordinates, FABL.tribesBlockEmbedding}
\difficulty{1}
Let $w,s$ be natural numbers, let $i$ be one of $s$ block indices, and let $o$ be one of
$w$ positions. The position corresponding to $(i,o)$ belongs to the set of coordinates
of block $i$.
\end{theorem}

\begin{theorem}[Bribing a tribes block to negative one]
\lean{FABL.bribeFunction_tribes_block_neg_one}
\label{FABL.bribeFunction_tribes_block_neg_one}
\leanok
\uses{FABL.Sign, FABL.andFunction, FABL.bribeFunction, FABL.bribeInput, FABL.inputBlock, FABL.orFunction, FABL.tribes, FABL.tribesBlockCoordinates}
\difficulty{5}
Let $w,s$ be natural numbers and choose a block index $i$ among the $s$ blocks. For
every input, bribing all $w$ coordinates in block $i$ of the tribes function to $-1$
makes the function constantly equal to $-1$.
\end{theorem}

\begin{theorem}[Tribes influence toward negative one]
\lean{FABL.coalitionalInfluenceToward_tribes_block_neg_one}
\label{FABL.coalitionalInfluenceToward_tribes_block_neg_one}
\leanok
\uses{FABL.BooleanFunction, FABL.BooleanFunction.toReal, FABL.Sign, FABL.SignCube, FABL.bribeFunction_tribes_block_neg_one, FABL.coalitionalInfluenceToward, FABL.coalitionalInfluenceToward_neg_one_eq_mean_change, FABL.fourierCoeff_tribes_empty, FABL.mean, FABL.mean_eq_fourierCoeff_empty, FABL.tribes, FABL.tribesBlockCoordinates, FABL.tribes_monotone}
\difficulty{4}
Let $w,s$ be nonnegative integers and let $i$ be one of the $s$ blocks. For the tribes
function on $s$ blocks of $w$ coordinates, the coalitional influence of block $i$ toward
$-1$ is $(1-2^{-w})^s$.
\end{theorem}

\begin{theorem}[Exact Maximum Coalitional Influence]
\lean{FABL.criticalTribes_towardTrue_max_exact}
\label{FABL.criticalTribes_towardTrue_max_exact}
\leanok
\uses{FABL.card_tribesBlockCoordinates, FABL.coalitionalInfluenceToward, FABL.coalitionalInfluenceToward_le_complement_probability, FABL.coalitionalInfluenceToward_tribes_block_neg_one, FABL.tribes, FABL.tribesBlockCoordinates, FABL.tribesCritical, FABL.tribesCriticalDimension, FABL.tribesCriticalProbabilityDeficit, FABL.tribesCriticalProbabilityDeficit_eq_pow, FABL.tribesCriticalSize, FABL.tribesCriticalSize_pos, FABL.tribesCritical_neg_one_probability_eq_half_sub_deficit}
\difficulty{5}
Let $w>0$, and consider the critical tribes Boolean function on its critical dimension,
with target value $-1$. Every set $J$ of at most $w$ coordinates has coalitional
influence toward $-1$ at most $1/2+d_w$, where $d_w$ is the function's probability
deficit from $1/2$; moreover, there is a set of exactly $w$ coordinates attaining this
bound.
\end{theorem}

\begin{theorem}[Lower bound on critical tribe size]
\lean{FABL.two_pow_sub_two_le_tribesCriticalSize}
\label{FABL.two_pow_sub_two_le_tribesCriticalSize}
\leanok
\uses{FABL.tribesCriticalSize, FABL.tribesCriticalSizeError_mem_Icc, FABL.tribesCriticalSize_eq_main_sub_error}
\difficulty{5}
For each integer $w\ge4$, define a candidate size $s$ by $1-(1-2^{-w})^s\le 1/2$, and
let the critical size be the greatest such candidate among integers at most $2^{w+2}$.
Then $2^{w-2}$ is at most the critical size.
\end{theorem}

\begin{theorem}[Critical dimension lower bound]
\lean{FABL.width_le_clog_tribesCriticalDimension}
\label{FABL.width_le_clog_tribesCriticalDimension}
\leanok
\uses{FABL.tribesCriticalDimension, FABL.tribesCriticalSize, FABL.two_pow_sub_two_le_tribesCriticalSize}
\difficulty{4}
For every integer $w\ge 4$, the width $w$ is at most the base-$2$ ceiling logarithm of
the tribes critical dimension at width $w$.
\end{theorem}

\begin{theorem}[Exact Maximum Influence Under a Logarithmic Size Bound]
\lean{FABL.criticalTribes_towardTrue_max_clog_exact}
\label{FABL.criticalTribes_towardTrue_max_clog_exact}
\leanok
\uses{FABL.coalitionalInfluenceToward, FABL.coalitionalInfluenceToward_le_complement_probability, FABL.criticalTribes_towardTrue_max_exact, FABL.tribesCritical, FABL.tribesCriticalDimension, FABL.tribesCriticalProbabilityDeficit, FABL.tribesCritical_neg_one_probability_eq_half_sub_deficit, FABL.width_le_clog_tribesCriticalDimension}
\difficulty{4}
Let $w\ge 4$, and consider the critical tribes Boolean function on its critical
dimension, with target value $-1$. For every set $J$ of coordinates whose size is at
most the base-$2$ logarithm rounded up of that dimension, its coalitional influence
toward $-1$ is at most $1/2+d_w$, where $d_w$ is the function's probability deficit from
$1/2$; moreover, some such set has size exactly $w$ and attains this bound.
\end{theorem}

\begin{theorem}[Critical Tribes Probability Excess]
\lean{FABL.criticalTribes_towardTrue_excess_isBigO}
\label{FABL.criticalTribes_towardTrue_excess_isBigO}
\leanok
\uses{FABL.tribesCriticalDimension, FABL.tribesCriticalProbabilityDeficit, FABL.tribesCriticalProbabilityDeficit_isBigO_log_dimension_div_dimension}
\difficulty{1}
For each block width $w$, let $N_w$ be the dimension of the critical tribes function and
let $d_w$ be its probability deficit from $1/2$. Then the excess $(1/2+d_w)-1/2$ is
$O\!\left(\log(N_w)/N_w\right)$ as $w\to\infty$.
\end{theorem}

\begin{theorem}[Equivalent Asymptotic Bounds for the Probability Deficit]
\lean{FABL.criticalTribes_towardTrue_omega_iff_deficit_omega}
\label{FABL.criticalTribes_towardTrue_omega_iff_deficit_omega}
\leanok
\uses{FABL.tribesCriticalDimension, FABL.tribesCriticalProbabilityDeficit}
\difficulty{1}
For each block width $w$, let $N_w$ be the critical tribes dimension and let $d_w$ be
its probability deficit from $1/2$. The function $\log(N_w)/N_w$ is $O(d_w)$ as
$w\to\infty$ if and only if it is $O((1/2+d_w)-1/2)$ as $w\to\infty$.
\end{theorem}

\begin{theorem}[Probability of negative signs on prescribed coordinates]
\lean{FABL.uniformProbability_all_neg_one_on_finset}
\label{FABL.uniformProbability_all_neg_one_on_finset}
\leanok
\uses{FABL.FixedSignCube, FABL.FreeSignCube, FABL.Sign, FABL.SignCube, FABL.signCubeSplitEquiv, FABL.uniformProbability}
\difficulty{5}
Let $n$ be a natural number and let $J$ be a set of coordinates in $\{1,\ldots,n\}$. For
a uniformly chosen sign vector in $\{-1,1\}^n$, the probability that every coordinate in
$J$ equals $-1$ is $2^{-|J|}$.
\end{theorem}

\begin{definition}[Remaining coordinates in a tribes block]
\lean{FABL.tribesRestCoordinates}
\label{FABL.tribesRestCoordinates}
\leanok
\uses{FABL.tribesBlockCoordinates, FABL.tribesCoord}
For natural numbers $w$ and $s$ and a coordinate $i \text{ in } \text{Fin}(sw)$, the
remaining coordinates are the coordinates in the block containing $i$, with $i$ itself
removed.
\end{definition}

\begin{theorem}[Membership in the corresponding block]
\lean{FABL.mem_tribesBlockCoordinates_self}
\label{FABL.mem_tribesBlockCoordinates_self}
\leanok
\uses{FABL.mem_tribesBlockCoordinates, FABL.tribesBlockCoordinates, FABL.tribesCoord}
\difficulty{3}
For natural numbers $w,s$ and an index $i$ among the $sw$ coordinates, let $(a,b)$ be
the pair of block and within-block indices corresponding to $i$ under the standard
identification of $\text{Fin}(sw)$ with $\text{Fin}(s)\times\text{Fin}(w)$. Then $i$
belongs to the set of coordinates in block $a$.
\end{theorem}

\begin{theorem}[Cardinality after removing one block coordinate]
\lean{FABL.card_tribesRestCoordinates}
\label{FABL.card_tribesRestCoordinates}
\leanok
\uses{FABL.card_tribesBlockCoordinates, FABL.mem_tribesBlockCoordinates_self, FABL.tribesRestCoordinates}
\difficulty{3}
Let $w$ and $s$ be natural numbers, and let $i$ be one of the $sw$ positions in $s$
blocks of $w$ positions each. The block containing $i$, with $i$ removed, has
cardinality $w-1$.
\end{theorem}

\begin{theorem}[Rest coordinates characterize block condition]
\lean{FABL.tribesRestTrue_iff_all_neg_on_restCoordinates}
\label{FABL.tribesRestTrue_iff_all_neg_on_restCoordinates}
\leanok
\uses{FABL.SignCube, FABL.TribesRestTrue, FABL.finProdFinEquiv_tribesCoord, FABL.mem_tribesBlockCoordinates, FABL.tribesBlockCoordinates, FABL.tribesCoord, FABL.tribesRestCoordinates}
\difficulty{5}
Let $w,s$ be natural numbers, let $i$ be a coordinate in the sign cube of dimension
$sw$, and let $x$ be an element of that cube. The condition that every coordinate of the
block containing $i$, other than the coordinate corresponding to $i$, has sign $-1$
holds if and only if every coordinate in the block's rest-coordinate set has sign $-1$.
\end{theorem}

\begin{theorem}[Probability of a tribes block condition]
\lean{FABL.uniformProbability_tribesRestTrue}
\label{FABL.uniformProbability_tribesRestTrue}
\leanok
\uses{FABL.SignCube, FABL.TribesRestTrue, FABL.card_tribesRestCoordinates, FABL.tribesRestCoordinates, FABL.tribesRestTrue_iff_all_neg_on_restCoordinates, FABL.uniformProbability, FABL.uniformProbability_all_neg_one_on_finset}
\difficulty{4}
For natural numbers $w,s$ and an index $i$ among the $sw$ coordinates, choose an input
uniformly from $\{−1,1\}^{sw}$ and view it as $s$ blocks of size $w$. The probability
that every coordinate in the block containing $i$, except the coordinate $i$ itself,
equals $−1$ is $2^{-(w-1)}$.
\end{theorem}

\begin{theorem}[Bribing commutes with an outside coordinate flip]
\lean{FABL.bribeInput_flipCoordinate_of_not_mem}
\label{FABL.bribeInput_flipCoordinate_of_not_mem}
\leanok
\uses{FABL.Sign, FABL.SignCube, FABL.bribeInput, FABL.flipCoordinate}
\difficulty{4}
Let $n$ be a natural number, let $J$ be a set of coordinates, choose a coordinate
$i\notin J$, a sign $b$, and an input $x$ in the $n$-dimensional sign cube. Flipping
coordinate $i$ before bribing the coordinates in $J$ to $b$ gives the same input as
bribing first and then flipping coordinate $i$.
\end{theorem}

\begin{theorem}[Pivotality in bribed tribes]
\lean{FABL.isPivotal_bribeFunction_tribes_imp_restTrue}
\label{FABL.isPivotal_bribeFunction_tribes_imp_restTrue}
\leanok
\uses{FABL.IsPivotal, FABL.SignCube, FABL.TribesRestTrue, FABL.bribeFunction, FABL.bribeInput, FABL.bribeInput_flipCoordinate_of_not_mem, FABL.finProdFinEquiv_tribesCoord, FABL.flipCoordinate, FABL.isPivotal_tribes_iff, FABL.tribes, FABL.tribesCoord}
\difficulty{5}
Let $w,s\in\mathbb N$, let $J\subseteq\{1,\ldots,sw\}$ be a set of coordinates, and let
$i$ be a coordinate outside $J$. For an input $x\in\{\pm1\}^{sw}$, if the tribes
function remains pivotal at $i$ after every coordinate in $J$ is set to $+1$, then every
other coordinate in the same block as $i$ has value $-1$ in $x$.
\end{theorem}

\begin{theorem}[Influence bound for bribed tribes]
\lean{FABL.booleanInfluence_bribeFunction_tribes_le}
\label{FABL.booleanInfluence_bribeFunction_tribes_le}
\leanok
\uses{FABL.IsPivotal, FABL.TribesRestTrue, FABL.booleanInfluence, FABL.booleanInfluence_bribeFunction_eq_zero_of_mem, FABL.bribeFunction, FABL.isPivotal_bribeFunction_tribes_imp_restTrue, FABL.tribes, FABL.uniformProbability, FABL.uniformProbability_mono, FABL.uniformProbability_tribesRestTrue}
\difficulty{4}
For nonnegative integers $w$ and $s$, a set $J$ of coordinates of the $(sw)$-dimensional
sign cube, and a coordinate $i$, consider the tribes function with $s$ blocks of width
$w$, after setting the coordinates in $J$ to $+1$. Its influence in coordinate $i$ is at
most $1/2^{w-1}$.
\end{theorem}

\begin{theorem}[Reciprocal powers of two]
\lean{FABL.two_mul_inv_two_pow_eq_inv_pred_pow}
\label{FABL.two_mul_inv_two_pow_eq_inv_pred_pow}
\leanok
\difficulty{3}
For every positive natural number $w$, $2\cdot(2^w)^{-1}=(2^{w-1})^{-1}$.
\end{theorem}

\begin{theorem}[Mean gain from bribing tribes]
\lean{FABL.mean_bribeFunction_tribes_one_sub_mean_le}
\label{FABL.mean_bribeFunction_tribes_one_sub_mean_le}
\leanok
\uses{FABL.BooleanFunction.toReal, FABL.booleanInfluence_bribeFunction_tribes_le, FABL.bribeFunction, FABL.bribeFunction_empty, FABL.mean, FABL.mean_bribeFunction_insert_one, FABL.tribes, FABL.tribes_monotone, FABL.two_mul_inv_two_pow_eq_inv_pred_pow}
\difficulty{5}
Let $w>0$ and $s\text{ be nonnegative integers}$, and let $J$ be a set of coordinates in
the sign cube $\text{of dimension }sw$. For the tribes function that applies AND within
each block of $w$ coordinates and OR across the $s$ blocks, forcing every coordinate in
$J$ to $+1$ increases its mean by at most $2|J|/2^w$; that is, the mean after forcing is
at most the original mean plus $2|J|/2^w$.
\end{theorem}

\begin{theorem}[Tribes influence upper bound]
\lean{FABL.coalitionalInfluenceToward_tribes_one_le_card_mul_inv_pow}
\label{FABL.coalitionalInfluenceToward_tribes_one_le_card_mul_inv_pow}
\leanok
\uses{FABL.coalitionalInfluenceToward, FABL.coalitionalInfluenceToward_one_eq_mean_change, FABL.mean_bribeFunction_tribes_one_sub_mean_le, FABL.tribes, FABL.tribes_monotone}
\difficulty{3}
Let $w,s$ be natural numbers with $w>0$, and let $J$ be any set of coordinates in the
sign cube of dimension $sw$. For the tribes function with $s$ blocks of width $w$, its
coalitional influence toward $1$ on $J$ is at most $|J|2^{-w}$.
\end{theorem}

\begin{theorem}[Influence Bound for Critical Tribes]
\lean{FABL.coalitionalInfluenceToward_criticalTribes_one_le}
\label{FABL.coalitionalInfluenceToward_criticalTribes_one_le}
\leanok
\uses{FABL.coalitionalInfluenceToward, FABL.coalitionalInfluenceToward_tribes_one_le_card_mul_inv_pow, FABL.tribesCritical, FABL.tribesCriticalDimension, FABL.tribesCriticalSize}
\difficulty{3}
For positive integers $w$ and any nonnegative integer $k$, let $s$ be the largest
integer at most $2^{w+2}$ such that $1-(1-2^{-w})^s\le 1/2$. Consider the Boolean
function on $sw$ sign coordinates formed by partitioning the coordinates into $s$ blocks
of width $w$, assigning each block the value $-1$ exactly when all its coordinates are
$-1$ (and $1$ otherwise), and assigning the full input the value $1$ exactly when every
block value is $1$ (and $-1$ otherwise). For every coalition $J$ of at most $k$
coordinates, its influence toward $1$—the fraction of inputs from which it can attain
$1$ minus the fraction already having value $1$—is at most $k2^{-w}$.
\end{theorem}

\begin{theorem}[Critical power relative error tends to zero]
\lean{FABL.tendsto_tribesCriticalPowerRelativeError_zero}
\label{FABL.tendsto_tribesCriticalPowerRelativeError_zero}
\leanok
\uses{FABL.tribesCriticalPowerRelativeError, FABL.tribesCriticalPowerRelativeError_isLittleO_one}
\difficulty{2}
As the positive integer $w$ tends to infinity, the quantity $2^w\log(D_w)/D_w-1$ tends
to $0$, where $D_w$ is the tribes critical dimension at $w$.
\end{theorem}

\begin{theorem}[Tribes dimension logarithmic bound]
\lean{FABL.eventually_inv_two_pow_le_two_mul_log_dimension_div_dimension}
\label{FABL.eventually_inv_two_pow_le_two_mul_log_dimension_div_dimension}
\leanok
\uses{FABL.eventually_two_pow_eq_dimension_div_log_mul_one_add_error, FABL.tendsto_tribesCriticalDimension_atTop, FABL.tendsto_tribesCriticalPowerRelativeError_zero, FABL.tribesCriticalDimension, FABL.tribesCriticalPowerRelativeError}
\difficulty{5}
For all sufficiently large integers $w$, $2^{-w}\leq 2\log(D_w)/D_w$, where $D_w$ is the
critical dimension defined from the largest critical tribes size up to $2^{w+2}$.
\end{theorem}

\begin{theorem}[Critical tribes coalitional influence bound]
\lean{FABL.eventually_coalitionalInfluenceToward_criticalTribes_one_le_log_bound}
\label{FABL.eventually_coalitionalInfluenceToward_criticalTribes_one_le_log_bound}
\leanok
\uses{FABL.coalitionalInfluenceToward, FABL.coalitionalInfluenceToward_criticalTribes_one_le, FABL.eventually_inv_two_pow_le_two_mul_log_dimension_div_dimension, FABL.tribesCritical, FABL.tribesCriticalDimension}
\difficulty{3}
For all sufficiently large natural numbers $w$, for every natural number $k$ and every
coalition $J$ of coordinates of the critical tribes function, if $|J|\le k$, then its
influence toward the value $1$ is at most $2k\log d/d$, where $d$ is the dimension of
that critical tribes function.
\end{theorem}

\begin{theorem}[Critical tribes influence bound]
\lean{FABL.coalitionalInfluence_criticalTribes_le_half_add_succ_card_mul_inv_pow}
\label{FABL.coalitionalInfluence_criticalTribes_le_half_add_succ_card_mul_inv_pow}
\leanok
\uses{FABL.coalitionalInfluence, FABL.coalitionalInfluenceToward_criticalTribes_one_le, FABL.coalitionalInfluenceToward_le_complement_probability, FABL.coalitionalInfluence_eq_toward_add_toward, FABL.tribesCritical, FABL.tribesCriticalDimension, FABL.tribesCriticalProbabilityDeficit_mem_Icc, FABL.tribesCritical_neg_one_probability_eq_half_sub_deficit}
\difficulty{4}
For a natural number $w>0$, let $s_w$ be the largest natural number $s\le 2^{w+2}$ such
that $1-(1-2^{-w})^s\le\tfrac12$, and take the tribes function with $s_w$ blocks of
width $w$. For every set $J$ of coordinates in its dimension $s_ww$, its coalitional
influence is at most $\tfrac12+(|J|+1)2^{-w}$.
\end{theorem}

\begin{definition}[Critical tribes coalition constant]
\lean{FABL.criticalTribesCoalitionConstant}
\label{FABL.criticalTribesCoalitionConstant}
\leanok
The critical tribes coalition constant is the real number $1/800$.
\end{definition}

\begin{theorem}[Positivity of the coalition constant]
\lean{FABL.criticalTribesCoalitionConstant_pos}
\label{FABL.criticalTribesCoalitionConstant_pos}
\leanok
\uses{FABL.criticalTribesCoalitionConstant}
\difficulty{1}
The critical tribes coalition constant, defined to be $\frac{1}{800}$, is positive.
\end{theorem}

\begin{theorem}[Logarithmic growth of critical dimension]
\lean{FABL.tendsto_log_tribesCriticalDimension_div_dimension_zero}
\label{FABL.tendsto_log_tribesCriticalDimension_div_dimension_zero}
\leanok
\uses{FABL.tendsto_tribesCriticalDimension_atTop, FABL.tribesCriticalDimension}
\difficulty{3}
For integers $w,s\geq0$, call $s$ a candidate at width $w$ if $1-(1-2^{-w})^s\leq 1/2$.
Let $s_w$ be the greatest candidate among the integers from $0$ through $2^{w+2}$,
taking $s_w=0$ if there is no candidate, and put $D_w=s_w w$. As $w\to+\infty$ through
integers, $\frac{\log D_w}{D_w}$ tends to $0$.
\end{theorem}

\begin{theorem}[Critical tribes coalitional influence bound]
\lean{FABL.eventually_criticalTribes_coalitionalInfluence_le_fiftyOnePercent}
\label{FABL.eventually_criticalTribes_coalitionalInfluence_le_fiftyOnePercent}
\leanok
\uses{FABL.coalitionalInfluence, FABL.coalitionalInfluence_criticalTribes_le_half_add_succ_card_mul_inv_pow, FABL.criticalTribesCoalitionConstant, FABL.criticalTribesCoalitionConstant_pos, FABL.eventually_inv_two_pow_le_two_mul_log_dimension_div_dimension, FABL.tendsto_log_tribesCriticalDimension_div_dimension_zero, FABL.tendsto_tribesCriticalDimension_atTop, FABL.tribesCritical, FABL.tribesCriticalDimension}
\difficulty{5}
For all sufficiently large natural numbers $w$, every coalition $J$ of coordinates of
the critical tribes function whose size is at most $d/(800\log d)$ has coalitional
influence at most $51/100$, where $d$ is the dimension of that function.
\end{theorem}

\begin{definition}[Majority coalition split equivalence]
\lean{FABL.majorityCoalitionSplitEquiv}
\label{FABL.majorityCoalitionSplitEquiv}
\leanok
\uses{FABL.SignCube}
For natural numbers $k$ and $m$, there is a bijection between sign assignments on $k+m$
coordinates and pairs consisting of a sign assignment on $k$ coordinates and one on $m$
coordinates. The bijection splits an assignment into its first $k$ and last $m$
coordinates, and its inverse concatenates the two assignments.
\end{definition}

\begin{definition}[Embedding into the first coordinates]
\lean{FABL.majorityHeadEmbedding}
\label{FABL.majorityHeadEmbedding}
\leanok
For natural numbers $k$ and $m$, the first $k$ coordinates embed into the coordinates
indexed by $0,\ldots,k+m-1$, by sending each index to the corresponding index.
\end{definition}

\begin{definition}[Majority head coalition]
\lean{FABL.majorityHeadCoalition}
\label{FABL.majorityHeadCoalition}
\leanok
\uses{FABL.majorityHeadEmbedding}
For natural numbers $k$ and $m$, the majority-head coalition is the set of coordinate
positions obtained by applying the majority-head embedding to every position among the
$k+m$ coordinates.
\end{definition}

\begin{theorem}[Cardinality of the majority head coalition]
\lean{FABL.card_majorityHeadCoalition}
\label{FABL.card_majorityHeadCoalition}
\leanok
\uses{FABL.majorityHeadCoalition}
\difficulty{1}
For natural numbers $k$ and $m$, the set of the first $k$ coordinates in $\bbr^{k+m}$ has
cardinality $k$.
\end{theorem}

\begin{theorem}[Membership of a left-summand index]
\lean{FABL.mem_majorityHeadCoalition_left}
\label{FABL.mem_majorityHeadCoalition_left}
\leanok
\uses{FABL.majorityHeadCoalition, FABL.majorityHeadEmbedding}
\difficulty{1}
For natural numbers $k,m$ and $i\in\{0,\ldots,k-1\}$, the index corresponding to
$(i,\mathrm{left})$ under the canonical identification of the disjoint union of sets of
sizes $k$ and $m$ with a set of size $k+m$ belongs to the coalition of the first $k$
indices.
\end{theorem}

\begin{theorem}[Right summand outside the head coalition]
\lean{FABL.not_mem_majorityHeadCoalition_right}
\label{FABL.not_mem_majorityHeadCoalition_right}
\leanok
\uses{FABL.majorityHeadCoalition}
\difficulty{4}
For natural numbers $k,m$ and any index $j<m$, the index corresponding to the right
summand element $j$ in the disjoint union of $k$ and $m$ indices does not belong to the
coalition consisting of the first $k$ indices.
\end{theorem}

\begin{theorem}[Membership in the head coalition]
\lean{FABL.mem_majorityHeadCoalition_castAdd}
\label{FABL.mem_majorityHeadCoalition_castAdd}
\leanok
\uses{FABL.majorityHeadCoalition, FABL.mem_majorityHeadCoalition_left}
\difficulty{3}
For natural numbers $k,m$ and an index $i\in\{0,\ldots,k-1\}$, the index obtained by
placing $i$ in the first summand of the disjoint union of sets of sizes $k$ and $m$
belongs to the coalition of the first $k$ indices in a set of size $k+m$.
\end{theorem}

\begin{theorem}[Tail indices outside the head coalition]
\lean{FABL.not_mem_majorityHeadCoalition_natAdd}
\label{FABL.not_mem_majorityHeadCoalition_natAdd}
\leanok
\uses{FABL.majorityHeadCoalition, FABL.not_mem_majorityHeadCoalition_right}
\difficulty{3}
For natural numbers $k,m$ and any index $j<m$, the index $k+j$ in the ordered set of
$k+m$ indices does not belong to the coalition consisting of the first $k$ indices.
\end{theorem}

\begin{theorem}[Expectation over a split sign cube]
\lean{FABL.expect_majorityCoalitionSplit_tail}
\label{FABL.expect_majorityCoalitionSplit_tail}
\leanok
\uses{FABL.SignCube, FABL.majorityCoalitionSplitEquiv}
\difficulty{4}
Let $k,m$ be nonnegative integers and let $g:\\{-1,1\\}^m\\to\\mathbb R$. For a
uniformly random sign vector $x\\in\\{-1,1\\}^{k+m}$, take its last $m$ coordinates,
denoted $y$. Then $\\mathbb E[g(y)]=\\mathbb E_{y\\in\\{-1,1\\}^m}[g(y)]$, where the
expectation on the right is uniform.
\end{theorem}

\begin{definition}[Majority after a head bribe]
\lean{FABL.majorityAfterHeadBribe}
\label{FABL.majorityAfterHeadBribe}
\leanok
\uses{FABL.BooleanFunction, FABL.Sign, FABL.signValue, FABL.thresholdSign}
Let $k,m$ be natural numbers and let $b\in\{ -1,1\}$. The Boolean function on $m$ sign
coordinates sends $y\in\{-1,1\}^{m}$ to the sign obtained by thresholding $k
b+\sum_{j=0}^{m-1}y_j$.
\end{definition}

\begin{theorem}[Majority after bribing the head coordinates]
\lean{FABL.bribeFunction_majority_head_apply}
\label{FABL.bribeFunction_majority_head_apply}
\leanok
\uses{FABL.Sign, FABL.SignCube, FABL.bribeFunction, FABL.bribeInput, FABL.majority, FABL.majorityAfterHeadBribe, FABL.majorityCoalitionSplitEquiv, FABL.majorityHeadCoalition, FABL.signValue}
\difficulty{4}
Let $k,m$ be natural numbers, let $b\in\{-1,1\}$, and let $x$ be a sign vector of length
$k+m$. Bribing the first $k$ coordinates of majority to $b$ yields the threshold sign of
$kb$ plus the sum of the last $m$ coordinates of $x$.
\end{theorem}

\begin{theorem}[Majority mean after bribing the first coordinates]
\lean{FABL.mean_bribeFunction_majority_head}
\label{FABL.mean_bribeFunction_majority_head}
\leanok
\uses{FABL.BooleanFunction.toReal, FABL.Sign, FABL.SignCube, FABL.bribeFunction, FABL.bribeFunction_majority_head_apply, FABL.expect_majorityCoalitionSplit_tail, FABL.majority, FABL.majorityAfterHeadBribe, FABL.majorityCoalitionSplitEquiv, FABL.majorityHeadCoalition, FABL.mean}
\difficulty{3}
For natural numbers $k,m$ and a sign $b$, consider majority on $k+m$ coordinates, with
the first $k$ coordinates replaced by $b$. Its mean equals the mean of the function on
the remaining $m$ coordinates that takes the threshold sign of $k b$ plus the sum of
those coordinates.
\end{theorem}

\begin{theorem}[Threshold sign of a natural number plus sign sum]
\lean{FABL.thresholdSign_nat_add_sign_sum_eq_one_iff}
\label{FABL.thresholdSign_nat_add_sign_sum_eq_one_iff}
\leanok
\uses{FABL.Sign, FABL.SignCube, FABL.signValue, FABL.thresholdSign, FABL.thresholdSign_of_neg, FABL.thresholdSign_of_nonneg}
\difficulty{4}
For natural numbers $k,m$ and a sign vector $y\\in\\{-1,1\\}^m$, the threshold sign of
$k+\\sum_{j=0}^{m-1}y_j$ is $1$ if and only if $-(k+1)<\\sum_{j=0}^{m-1}y_j$, where the
sum on the right is taken in the integers.
\end{theorem}

\begin{theorem}[Threshold sign of a signed sum]
\lean{FABL.thresholdSign_sign_sum_sub_nat_eq_one_iff}
\label{FABL.thresholdSign_sign_sum_sub_nat_eq_one_iff}
\leanok
\uses{FABL.Sign, FABL.SignCube, FABL.signValue, FABL.thresholdSign, FABL.thresholdSign_of_neg, FABL.thresholdSign_of_nonneg}
\difficulty{4}
Let $k,m$ be natural numbers, and let $y$ assign a sign in $\{-1,1\}$ to each of the $m$
coordinates. The threshold sign of the sum of these signs minus $k$ is $1$ if and only
if the sum of their integer values is strictly greater than $k-1$.
\end{theorem}

\begin{definition}[Majority plus bribe threshold]
\lean{FABL.majorityPlusBribeThreshold}
\label{FABL.majorityPlusBribeThreshold}
\leanok
For natural numbers $k$ and $m$, define the majority-plus-bribe threshold by
$-\frac{k+1}{\sqrt{m}}$.
\end{definition}

\begin{definition}[Majority minus bribe threshold]
\lean{FABL.majorityMinusBribeThreshold}
\label{FABL.majorityMinusBribeThreshold}
\leanok
For natural numbers $k$ and $m$, define the majority-minus-bribe threshold to be
$\frac{k-1}{\sqrt{m}}$.
\end{definition}

\begin{theorem}[Positive Bribe Tail Indicator]
\lean{FABL.majorityAfterHeadBribe_one_indicator}
\label{FABL.majorityAfterHeadBribe_one_indicator}
\leanok
\uses{FABL.SignCube, FABL.hammingUpperTailIndicator, FABL.linearForm, FABL.majorityAfterHeadBribe, FABL.majorityPlusBribeThreshold, FABL.normalizedRademacherSum, FABL.signValue, FABL.signValue_one, FABL.thresholdSign_nat_add_sign_sum_eq_one_iff}
\difficulty{4}
Let $k,m$ be nonnegative integers with $m>0$, and let $y\in\{−1,1\}^m$. The indicator
that the majority function with bribe $1$ returns $1$ at $y$ equals the indicator that
$m^{-1/2}\sum_{j=1}^m y_j>-(k+1)/\sqrt m$.
\end{theorem}

\begin{theorem}[Negative Bribe Tail Indicator]
\lean{FABL.majorityAfterHeadBribe_neg_one_one_indicator}
\label{FABL.majorityAfterHeadBribe_neg_one_one_indicator}
\leanok
\uses{FABL.SignCube, FABL.hammingUpperTailIndicator, FABL.linearForm, FABL.majorityAfterHeadBribe, FABL.majorityMinusBribeThreshold, FABL.normalizedRademacherSum, FABL.signValue, FABL.signValue_neg_one, FABL.thresholdSign_sign_sum_sub_nat_eq_one_iff}
\difficulty{4}
Let $k,m$ be nonnegative integers with $m>0$, and let $y\in\{−1,1\}^m$. The indicator
that the majority function with bribe $-1$ returns $1$ at $y$ equals the indicator that
$m^{-1/2}\sum_{j=1}^m y_j>(k-1)/\sqrt m$.
\end{theorem}

\begin{theorem}[Bribed majority bias identity]
\lean{FABL.positiveBias_bribeFunction_majority_head}
\label{FABL.positiveBias_bribeFunction_majority_head}
\leanok
\uses{FABL.Sign, FABL.SignCube, FABL.bribeFunction, FABL.bribeFunction_majority_head_apply, FABL.expect_majorityCoalitionSplit_tail, FABL.majority, FABL.majorityAfterHeadBribe, FABL.majorityCoalitionSplitEquiv, FABL.majorityHeadCoalition, FABL.positiveBias, FABL.uniformProbability}
\difficulty{3}
For any nonnegative integers $k,m$ and any sign $b\in\{-1,1\}$, the fraction of inputs
in the $(k+m)$-dimensional sign cube on which majority is positive after the first $k$
coordinates are set to $b$ equals the fraction of inputs $y$ in the $m$-dimensional sign
cube for which the threshold function with value $+1$ when $k b+\sum_{j=1}^m y_j\ge 0$
is positive.
\end{theorem}

\begin{theorem}[Positive bribery majority tail identity]
\lean{FABL.positiveBias_bribeFunction_majority_head_one}
\label{FABL.positiveBias_bribeFunction_majority_head_one}
\leanok
\uses{FABL.SignCube, FABL.bribeFunction, FABL.hammingUpperTailIndicator, FABL.majority, FABL.majorityAfterHeadBribe_one_indicator, FABL.majorityHeadCoalition, FABL.majorityPlusBribeThreshold, FABL.positiveBias, FABL.positiveBias_bribeFunction_majority_head, FABL.uniformProbability}
\difficulty{3}
Let $k$ be a nonnegative integer, let $m$ be a positive integer, and set the first $k$
coordinates of majority on the $(k+m)$-dimensional sign cube to $+1$. Its positive bias
equals the average over all $m$-dimensional sign vectors $y$ of the indicator that
$\sum_{j=1}^m y_j/\sqrt m>-(k+1)/\sqrt m$.
\end{theorem}

\begin{theorem}[Negative bribery majority tail identity]
\lean{FABL.positiveBias_bribeFunction_majority_head_neg_one}
\label{FABL.positiveBias_bribeFunction_majority_head_neg_one}
\leanok
\uses{FABL.SignCube, FABL.bribeFunction, FABL.hammingUpperTailIndicator, FABL.majority, FABL.majorityAfterHeadBribe_neg_one_one_indicator, FABL.majorityHeadCoalition, FABL.majorityMinusBribeThreshold, FABL.positiveBias, FABL.positiveBias_bribeFunction_majority_head, FABL.uniformProbability}
\difficulty{3}
Let $k$ be a nonnegative integer, let $m$ be a positive integer, and set the first $k$
coordinates of majority on the $(k+m)$-dimensional sign cube to $-1$. Its positive bias
equals the average over all $m$-dimensional sign vectors $y$ of the indicator that
$\sum_{j=1}^m y_j/\sqrt m>(k-1)/\sqrt m$.
\end{theorem}

\begin{theorem}[Majority bias as an upper tail]
\lean{FABL.positiveBias_majority_eq_hammingUpperTail}
\label{FABL.positiveBias_majority_eq_hammingUpperTail}
\leanok
\uses{FABL.SignCube, FABL.hammingUpperTailIndicator, FABL.majority, FABL.majorityAfterHeadBribe, FABL.majorityAfterHeadBribe_one_indicator, FABL.majorityPlusBribeThreshold, FABL.positiveBias, FABL.signValue_one, FABL.uniformProbability}
\difficulty{3}
For every positive integer $n$, the positive bias of majority on the $n$-dimensional
sign cube equals the fraction of sign vectors $y$ for which $\sum_{i=1}^n y_i/\sqrt
n>-1/\sqrt n$.
\end{theorem}

\begin{theorem}[Positive Influence as a Positive Bias Change]
\lean{FABL.coalitionalInfluenceToward_one_eq_positiveBias_change}
\label{FABL.coalitionalInfluenceToward_one_eq_positiveBias_change}
\leanok
\uses{FABL.BooleanFunction, FABL.bribeFunction, FABL.coalitionalInfluenceToward, FABL.coalitionalInfluenceToward_one_eq_mean_change, FABL.positiveBias, FABL.positiveBias_eq_mean}
\difficulty{2}
Let $n$ be a nonnegative integer, let $f:\\{-1,1\\}^n\\to\\{-1,1\\}$ be monotone in the
coordinatewise order, and let $J\\subseteq\\{0,\\ldots,n-1\\}$. Write $f_{J,+}$ for the
function obtained by setting every coordinate in $J$ to $+1$ before applying $f$, and
write the positive bias of a function for the uniform probability that its value is
$+1$. Then the coalitional influence of $f$ toward $+1$ equals the positive bias of
$f_{J,+}$ minus the positive bias of $f$.
\end{theorem}

\begin{theorem}[Negative Influence as a Positive Bias Change]
\lean{FABL.coalitionalInfluenceToward_neg_one_eq_positiveBias_change}
\label{FABL.coalitionalInfluenceToward_neg_one_eq_positiveBias_change}
\leanok
\uses{FABL.BooleanFunction, FABL.bribeFunction, FABL.coalitionalInfluenceToward, FABL.coalitionalInfluenceToward_neg_one_eq_mean_change, FABL.positiveBias, FABL.positiveBias_eq_mean}
\difficulty{2}
Let $n$ be a nonnegative integer, let $f:\\{-1,1\\}^n\\to\\{-1,1\\}$ be monotone in the
coordinatewise order, and let $J\\subseteq\\{0,\\ldots,n-1\\}$. Write $f_{J,-}$ for the
function obtained by setting every coordinate in $J$ to $-1$ before applying $f$, and
write the positive bias of a function for the uniform probability that its value is
$+1$. Then the coalitional influence of $f$ toward $-1$ equals the positive bias of $f$
minus the positive bias of $f_{J,-}$.
\end{theorem}

\begin{theorem}[Coalitional influence toward majority head one]
\lean{FABL.coalitionalInfluenceToward_majority_head_one}
\statementsource{Analysis-of-Boolean-Functions-by-Ryan-ODonnell}{pdf-ffbe9c7b3306-p275-b002}
\label{FABL.coalitionalInfluenceToward_majority_head_one}
\leanok
\uses{FABL.SignCube, FABL.coalitionalInfluenceToward, FABL.coalitionalInfluenceToward_one_eq_positiveBias_change, FABL.hammingUpperTailIndicator, FABL.majority, FABL.majorityHeadCoalition, FABL.majorityPlusBribeThreshold, FABL.majority_monotone, FABL.positiveBias_bribeFunction_majority_head_one, FABL.positiveBias_majority_eq_hammingUpperTail}
\difficulty{3}
% dataset-pipeline: stub -- prose pending build_informal_statements.py
Statement pending regeneration.
\end{theorem}

\begin{theorem}[Coalitional influence toward majority head neg one]
\lean{FABL.coalitionalInfluenceToward_majority_head_neg_one}
\statementsource{Analysis-of-Boolean-Functions-by-Ryan-ODonnell}{pdf-ffbe9c7b3306-p275-b002}
\label{FABL.coalitionalInfluenceToward_majority_head_neg_one}
\leanok
\uses{FABL.SignCube, FABL.coalitionalInfluenceToward, FABL.coalitionalInfluenceToward_neg_one_eq_positiveBias_change, FABL.hammingUpperTailIndicator, FABL.majority, FABL.majorityHeadCoalition, FABL.majorityMinusBribeThreshold, FABL.majority_monotone, FABL.positiveBias_bribeFunction_majority_head_neg_one, FABL.positiveBias_majority_eq_hammingUpperTail}
\difficulty{3}
% dataset-pipeline: stub -- prose pending build_informal_statements.py
Statement pending regeneration.
\end{theorem}

\begin{theorem}[Gaussian approximation for a Rademacher tail]
\lean{FABL.abs_expect_hammingUpperTailIndicator_dimension_le}
\label{FABL.abs_expect_hammingUpperTailIndicator_dimension_le}
\leanok
\uses{FABL.SignCube, FABL.abs_expect_hammingUpperTailIndicator_sub_standardGaussianUpperTail_le, FABL.hammingUpperTailIndicator, FABL.standardGaussianUpperTail}
\difficulty{4}
For every real threshold $u$ and positive integer $n$, the absolute difference between
the probability that the normalized Rademacher sum on the $n$-dimensional sign cube
exceeds $u$ and the standard Gaussian upper-tail probability at $u$ is at most $2C
n^{-1/2}$, where $C$ is the third-moment Berry–Esseen constant.
\end{theorem}

\begin{theorem}[Finite-sample error for majority influence]
\lean{FABL.majority_head_directed_influence_finite_error}
\label{FABL.majority_head_directed_influence_finite_error}
\leanok
\uses{FABL.SignCube, FABL.abs_expect_hammingUpperTailIndicator_dimension_le, FABL.coalitionalInfluenceToward, FABL.coalitionalInfluenceToward_majority_head_neg_one, FABL.coalitionalInfluenceToward_majority_head_one, FABL.hammingUpperTailIndicator, FABL.majority, FABL.majorityHeadCoalition, FABL.majorityMinusBribeThreshold, FABL.majorityPlusBribeThreshold, FABL.standardGaussianUpperTail}
\difficulty{5}
Let $k,m$ be natural numbers with $m>0$. For majority on $k+m$ signs and the coalition
consisting of the first $k$ coordinates, the influence toward $+1$ differs from the
difference between the standard Gaussian upper tails at $-\frac{k+1}{\sqrt m}$ and
$-\frac{1}{\sqrt{k+m}}$ by at most
\[
2C\left(\frac{1}{\sqrt m}+\frac{1}{\sqrt{k+m}}\right),
\]
and the influence toward $-1$ differs from the difference between the upper tails at
$-\frac{1}{\sqrt{k+m}}$ and $\frac{k-1}{\sqrt m}$ by at most the same quantity, where
$C$ is the third-moment Berry–Esseen constant.
\end{theorem}

\begin{theorem}[Gaussian limit for moving thresholds]
\lean{FABL.tendsto_expect_hammingUpperTailIndicator_of_tendsto}
\label{FABL.tendsto_expect_hammingUpperTailIndicator_of_tendsto}
\leanok
\uses{FABL.SignCube, FABL.abs_expect_hammingUpperTailIndicator_dimension_le, FABL.continuous_standardGaussianUpperTail, FABL.hammingUpperTailIndicator, FABL.standardGaussianUpperTail}
\difficulty{6}
Let $d:\mathbb N\to\mathbb N$ and $u:\mathbb N\to\mathbb R$ be sequences, and let
$a\in\mathbb R$. If $d(q)\to\infty$ and $u(q)\to a$ as $q\to\infty$, then the expected
indicator that the normalized Rademacher sum on a uniformly random sign vector of
dimension $d(q)$ exceeds $u(q)$ converges to the standard Gaussian upper tail at $a$.
\end{theorem}

\begin{definition}[Gaussian majority limit profile]
\lean{FABL.majorityCoalitionGaussianLimit}
\statementsource{Analysis-of-Boolean-Functions-by-Ryan-ODonnell}{
  pdf-ffbe9c7b3306-p118-b004,
  pdf-ffbe9c7b3306-p119-b001,
  pdf-ffbe9c7b3306-p120-b003,
  pdf-ffbe9c7b3306-p275-b002
}
\label{FABL.majorityCoalitionGaussianLimit}
\leanok
\uses{FABL.standardGaussianUpperTail}
For $t\in\mathbb{R}$, define the Gaussian majority limit profile to be $\frac12$ minus
the upper-tail probability $\Pr[Z>t]$, where $Z$ is a standard normal random variable.
\end{definition}

\begin{theorem}[Standard Gaussian upper tail at zero]
\lean{FABL.standardGaussianUpperTail_zero}
\label{FABL.standardGaussianUpperTail_zero}
\leanok
\uses{FABL.standardGaussianUpperTail, FABL.standardGaussianUpperTail_neg}
\difficulty{3}
The standard Gaussian upper-tail probability at $0$ is $1/2$.
\end{theorem}

\begin{theorem}[Limit of the majority plus-bribe threshold]
\lean{FABL.tendsto_majorityPlusBribeThreshold}
\label{FABL.tendsto_majorityPlusBribeThreshold}
\leanok
\uses{FABL.majorityPlusBribeThreshold}
\difficulty{4}
Let $k,m:\mathbb{N}\to\mathbb{N}$ and $t\in\mathbb{R}$. If $m(q)\to\infty$ and
$k(q)/\sqrt{m(q)}\to t$ as $q\to\infty$, then $-(k(q)+1)/\sqrt{m(q)}\to -t$.
\end{theorem}

\begin{theorem}[Limit of the majority minus-bribe threshold]
\lean{FABL.tendsto_majorityMinusBribeThreshold}
\label{FABL.tendsto_majorityMinusBribeThreshold}
\leanok
\uses{FABL.majorityMinusBribeThreshold}
\difficulty{4}
Let $k,m:\mathbb{N}\to\mathbb{N}$ and $t\in\mathbb{R}$. If $m(q)\to\infty$ and
$k(q)/\sqrt{m(q)}\to t$ as $q\to\infty$, then $\bigl(k(q)-1\bigr)/\sqrt{m(q)}\to t$.
\end{theorem}

\begin{theorem}[Majority baseline threshold limit]
\lean{FABL.tendsto_majorityBaselineThreshold}
\label{FABL.tendsto_majorityBaselineThreshold}
\leanok
\difficulty{4}
Let $k,m:\mathbb{N}\to\mathbb{N}$ be sequences, and suppose $m(q)\to+\infty$ as
$q\to+\infty$. Then $-1/\sqrt{k(q)+m(q)}$ tends to $0$ as $q\to+\infty$.
\end{theorem}

\begin{theorem}[Tendsto majority head directed influences]
\lean{FABL.tendsto_majority_head_directed_influences}
\statementsource{Analysis-of-Boolean-Functions-by-Ryan-ODonnell}{pdf-ffbe9c7b3306-p275-b002}
\label{FABL.tendsto_majority_head_directed_influences}
\leanok
\uses{FABL.coalitionalInfluenceToward, FABL.coalitionalInfluenceToward_majority_head_neg_one, FABL.coalitionalInfluenceToward_majority_head_one, FABL.majority, FABL.majorityCoalitionGaussianLimit, FABL.majorityHeadCoalition, FABL.majorityMinusBribeThreshold, FABL.majorityPlusBribeThreshold, FABL.standardGaussianUpperTail, FABL.standardGaussianUpperTail_neg, FABL.standardGaussianUpperTail_zero, FABL.tendsto_expect_hammingUpperTailIndicator_of_tendsto, FABL.tendsto_majorityBaselineThreshold, FABL.tendsto_majorityMinusBribeThreshold, FABL.tendsto_majorityPlusBribeThreshold}
\difficulty{5}
% dataset-pipeline: stub -- prose pending build_informal_statements.py
Statement pending regeneration.
\end{theorem}

\begin{theorem}[Total influence of the majority head coalition]
\lean{FABL.tendsto_majority_head_total_coalitionalInfluence}
\label{FABL.tendsto_majority_head_total_coalitionalInfluence}
\leanok
\uses{FABL.coalitionalInfluence, FABL.coalitionalInfluenceToward, FABL.coalitionalInfluence_eq_toward_add_toward, FABL.majority, FABL.majorityCoalitionGaussianLimit, FABL.majorityHeadCoalition, FABL.tendsto_majority_head_directed_influences}
\difficulty{4}
Let $k,m:\mathbb{N}\to\mathbb{N}$ and $t\in\mathbb{R}$. If $m(q)\to\infty$ and
$k(q)/\sqrt{m(q)}\to t$, then the probability that the first $k(q)$ coordinates can make
majority attain both possible outputs, for a uniformly random input on $k(q)+m(q)$
coordinates, converges to $2\bigl(\tfrac12-Q(t)\bigr)$, where $Q(t)=\Pr(Z>t)$ for a
standard normal random variable $Z$.
\end{theorem}

\begin{theorem}[Limit of the outside-coordinate ratio]
\lean{FABL.tendsto_majority_outside_ratio_of_total_ratio}
\label{FABL.tendsto_majority_outside_ratio_of_total_ratio}
\leanok
\difficulty{6}
Let $k,m:\bbn\to\bbn$ be sequences of natural numbers, and let $t\in\bbr$. If $m(q)\to\infty$
and $k(q)/\sqrt{k(q)+m(q)}\to t$, then $k(q)/\sqrt{m(q)}\to t$.
\end{theorem}

\begin{theorem}[Directed influence of the majority head coalition]
\lean{FABL.tendsto_majority_head_directed_influences_totalDimension}
\label{FABL.tendsto_majority_head_directed_influences_totalDimension}
\leanok
\uses{FABL.coalitionalInfluenceToward, FABL.majority, FABL.majorityCoalitionGaussianLimit, FABL.majorityHeadCoalition, FABL.tendsto_majority_head_directed_influences, FABL.tendsto_majority_outside_ratio_of_total_ratio}
\difficulty{1}
Let $k,m:\mathbb{N}\to\mathbb{N}$ and $t\in\mathbb{R}$. Suppose $m(q)\to\infty$ and
$k(q)/\sqrt{k(q)+m(q)}\to t$. For majority on $k(q)+m(q)$ coordinates, the head
coalition consisting of the first $k(q)$ coordinates has influence toward each of $+1$
and $-1$ converging to $\tfrac12-Q(t)$, where $Q(t)=\Pr(Z>t)$ for a standard normal
random variable $Z$.
\end{theorem}

\begin{theorem}[Total influence of the majority head coalition]
\lean{FABL.tendsto_majority_head_total_coalitionalInfluence_totalDimension}
\proofsource{Analysis-of-Boolean-Functions-by-Ryan-ODonnell}{pdf-ffbe9c7b3306-p275-b002}
\label{FABL.tendsto_majority_head_total_coalitionalInfluence_totalDimension}
\leanok
\uses{FABL.coalitionalInfluence, FABL.majority, FABL.majorityCoalitionGaussianLimit, FABL.majorityHeadCoalition, FABL.tendsto_majority_head_total_coalitionalInfluence, FABL.tendsto_majority_outside_ratio_of_total_ratio}
\difficulty{1}
Let $k,m:\mathbb{N}\to\mathbb{N}$ and $t\in\mathbb{R}$. If $m(q)\to\infty$ and
$k(q)/\sqrt{k(q)+m(q)}\to t$, then the probability that the first $k(q)$ coordinates can
make majority attain both possible outputs, for a uniformly random input on $k(q)+m(q)$
coordinates, converges to $2\bigl(\tfrac12-Q(t)\bigr)$, where $Q(t)=\Pr(Z>t)$ for a
standard normal random variable $Z$.
\end{theorem}

\begin{theorem}[Vanishing Coalitional Influence of a Majority Head]
\lean{FABL.tendsto_majority_head_total_coalitionalInfluence_zero_of_subsqrt}
\label{FABL.tendsto_majority_head_total_coalitionalInfluence_zero_of_subsqrt}
\leanok
\uses{FABL.coalitionalInfluence, FABL.majority, FABL.majorityCoalitionGaussianLimit, FABL.majorityHeadCoalition, FABL.standardGaussianUpperTail_zero, FABL.tendsto_majority_head_total_coalitionalInfluence_totalDimension}
\difficulty{3}
Let $k,m:\bbn\to\bbn$ be sequences of natural numbers. If $m(q)\to\infty$ and
$k(q)/\sqrt{k(q)+m(q)}\to 0$ as $q\to\infty$, then the coalitional influence of the
first $k(q)$ coordinates for the majority function on $k(q)+m(q)$ sign coordinates tends
to zero as $q\to\infty$.
\end{theorem}

\begin{theorem}[Tail probability convergence from Gaussian tails]
\lean{FABL.tendsto_expect_hammingUpperTailIndicator_of_tail_tendsto}
\label{FABL.tendsto_expect_hammingUpperTailIndicator_of_tail_tendsto}
\leanok
\uses{FABL.SignCube, FABL.abs_expect_hammingUpperTailIndicator_dimension_le, FABL.hammingUpperTailIndicator, FABL.standardGaussianUpperTail}
\difficulty{5}
Let $d:\mathbb N\to\mathbb N$ and $u:\mathbb N\to\mathbb R$ be sequences, and let
$L\in\mathbb R$. If $d(q)\to\infty$ and the standard Gaussian upper tail at $u(q)$
converges to $L$ as $q\to\infty$, then the expected indicator that the normalized
Rademacher sum on a uniformly random sign vector of dimension $d(q)$ exceeds $u(q)$
converges to $L$.
\end{theorem}

\begin{theorem}[Divergence of the outside-coordinate ratio]
\lean{FABL.tendsto_majority_outside_ratio_atTop_of_total_ratio_atTop}
\label{FABL.tendsto_majority_outside_ratio_atTop_of_total_ratio_atTop}
\leanok
\difficulty{4}
Let $k,m:\bbn\to\bbn$ be sequences of natural numbers. If $m(q)\to\infty$ and
$k(q)/\sqrt{k(q)+m(q)}\to\infty$, then $k(q)/\sqrt{m(q)}\to\infty$.
\end{theorem}

\begin{theorem}[Total influence when the head coalition dominates]
\lean{FABL.tendsto_majority_head_total_coalitionalInfluence_one_of_omega_sqrt}
\label{FABL.tendsto_majority_head_total_coalitionalInfluence_one_of_omega_sqrt}
\leanok
\uses{FABL.SignCube, FABL.coalitionalInfluence, FABL.coalitionalInfluenceToward, FABL.coalitionalInfluenceToward_majority_head_neg_one, FABL.coalitionalInfluenceToward_majority_head_one, FABL.coalitionalInfluence_eq_toward_add_toward, FABL.hammingUpperTailIndicator, FABL.majority, FABL.majorityHeadCoalition, FABL.majorityMinusBribeThreshold, FABL.majorityPlusBribeThreshold, FABL.standardGaussianUpperTail_zero, FABL.tendsto_expect_hammingUpperTailIndicator_of_tail_tendsto, FABL.tendsto_expect_hammingUpperTailIndicator_of_tendsto, FABL.tendsto_majorityBaselineThreshold, FABL.tendsto_majority_outside_ratio_atTop_of_total_ratio_atTop, FABL.tendsto_standardGaussianUpperTail_atBot, FABL.tendsto_standardGaussianUpperTail_atTop}
\difficulty{6}
Let $k,m:\mathbb{N}\to\mathbb{N}$. If $m(q)\to\infty$ and
$k(q)/\sqrt{k(q)+m(q)}\to\infty$, then the probability that the first $k(q)$ coordinates
can make majority attain both possible outputs, for a uniformly random input on
$k(q)+m(q)$ coordinates, converges to $1$.
\end{theorem}

\begin{definition}[Coalition transport permutation]
\lean{FABL.coalitionTransportPerm}
\label{FABL.coalitionTransportPerm}
\leanok
Let $n$ be a natural number, and let $A$ and $B$ be subsets of $\{0,\ldots,n-1\}$ with
equal cardinality. There is a permutation of $\{0,\ldots,n-1\}$ obtained by extending a
bijection between $A$ and $B$.
\end{definition}

\begin{theorem}[Coalition transport preserves membership]
\lean{FABL.mem_coalitionTransportPerm_iff}
\label{FABL.mem_coalitionTransportPerm_iff}
\leanok
\uses{FABL.coalitionTransportPerm}
\difficulty{4}
Let $A$ and $B$ be subsets of $\{0,\ldots,n-1\}$ with equal cardinality, and let
$i\in\{0,\ldots,n-1\}$. The image of $i$ under the coalition transport permutation
belongs to $B$ if and only if $i$ belongs to $A$.
\end{theorem}

\begin{theorem}[Bribing under coalition transport]
\lean{FABL.bribeInput_permuteInput_coalitionTransport}
\label{FABL.bribeInput_permuteInput_coalitionTransport}
\leanok
\uses{FABL.Sign, FABL.SignCube, FABL.bribeInput, FABL.coalitionTransportPerm, FABL.mem_coalitionTransportPerm_iff, FABL.permuteInput}
\difficulty{4}
Let $n$ be a nonnegative integer, and let $A$ and $B$ be sets of coordinates of equal
size. Choose a permutation of the $n$ coordinates that transports $A$ to $B$ by a
bijection on those sets and is extended to the remaining coordinates as in the transport
construction. For every sign $b$ and sign vector $x$, replacing the coordinates in $A$
by $b$ after permuting $x$ gives the same vector as permuting the vector obtained by
replacing the coordinates in $B$ by $b$.
\end{theorem}

\begin{theorem}[Majority mean depends on coalition size]
\lean{FABL.mean_bribeFunction_majority_eq_of_card}
\label{FABL.mean_bribeFunction_majority_eq_of_card}
\leanok
\uses{FABL.BooleanFunction.toReal, FABL.Sign, FABL.bribeFunction, FABL.bribeInput, FABL.bribeInput_permuteInput_coalitionTransport, FABL.coalitionTransportPerm, FABL.majority, FABL.majority_symmetric, FABL.mean, FABL.permuteInput, FABL.permuteInputEquiv, FABL.signValue}
\difficulty{4}
Let $A$ and $B$ be sets of coordinates in the $n$-dimensional sign cube with equal
cardinality, and let $b$ be a sign. The means of majority after replacing the
coordinates in $A$ by $b$ and after replacing the coordinates in $B$ by $b$ are equal.
\end{theorem}

\begin{theorem}[Majority Influence Depends on Coalition Size]
\lean{FABL.coalitionalInfluenceToward_majority_eq_of_card}
\label{FABL.coalitionalInfluenceToward_majority_eq_of_card}
\leanok
\uses{FABL.Sign, FABL.coalitionalInfluenceToward, FABL.coalitionalInfluenceToward_neg_one_eq_mean_change, FABL.coalitionalInfluenceToward_one_eq_mean_change, FABL.majority, FABL.majority_monotone, FABL.mean_bribeFunction_majority_eq_of_card}
\difficulty{4}
Let $n$ be a nonnegative integer, and let $A,B\\subseteq\\{0,\\ldots,n-1\\}$ be sets of
coordinates with $|A|=|B|$. On the sign cube $\\{-1,1\\}^n$, define majority to be $+1$
when the sum of the coordinates is nonnegative and $-1$ otherwise. For either target
sign $b\\in\\{-1,1\\}$, the coalitional influence of majority toward $b$ is the same for
$A$ and $B$.
\end{theorem}

\begin{theorem}[Majority influence depends on set size]
\lean{FABL.coalitionalInfluence_majority_eq_of_card}
\label{FABL.coalitionalInfluence_majority_eq_of_card}
\leanok
\uses{FABL.coalitionalInfluence, FABL.coalitionalInfluenceToward_majority_eq_of_card, FABL.coalitionalInfluence_eq_toward_add_toward, FABL.majority}
\difficulty{3}
Let $A$ and $B$ be sets of coordinates in an $n$-dimensional sign cube, and suppose
$|A|=|B|$. The coalitional influences of the majority function on $A$ and $B$ are equal.
\end{theorem}

\begin{theorem}[Tendsto majority arbitrary coalition total]
\lean{FABL.tendsto_majority_arbitraryCoalition_total}
\statementsource{Analysis-of-Boolean-Functions-by-Ryan-ODonnell}{pdf-ffbe9c7b3306-p275-b002}
\label{FABL.tendsto_majority_arbitraryCoalition_total}
\leanok
\uses{FABL.coalitionalInfluence, FABL.coalitionalInfluence_majority_eq_of_card, FABL.majority, FABL.majorityCoalitionGaussianLimit, FABL.majorityHeadCoalition, FABL.tendsto_majority_head_total_coalitionalInfluence_totalDimension}
\difficulty{3}
% dataset-pipeline: stub -- prose pending build_informal_statements.py
Statement pending regeneration.
\end{theorem}

\begin{theorem}[Coalitional influence of majority for arbitrary coalitions]
\lean{FABL.tendsto_majority_arbitraryCoalition_directed}
\proofsource{Analysis-of-Boolean-Functions-by-Ryan-ODonnell}{pdf-ffbe9c7b3306-p275-b002}
\label{FABL.tendsto_majority_arbitraryCoalition_directed}
\leanok
\uses{FABL.coalitionalInfluenceToward, FABL.coalitionalInfluenceToward_majority_eq_of_card, FABL.majority, FABL.majorityCoalitionGaussianLimit, FABL.majorityHeadCoalition, FABL.tendsto_majority_head_directed_influences_totalDimension}
\difficulty{4}
Let $k,m:\mathbb{N}\to\mathbb{N}$, and for each $q$ let $J_q$ be a set of coordinates
among $k(q)+m(q)$ coordinates, with $|J_q|=k(q)$. Suppose $m(q)\to\infty$ and
$k(q)/\sqrt{k(q)+m(q)}\to t\in\mathbb{R}$. For majority on the sign cube, the influence
of $J_q$ toward either output, defined as the probability that the coalition can attain
that output minus its probability under a uniformly random input, converges to
$\tfrac12-Q(t)$, where $Q(t)=\Pr(Z>t)$ for a standard normal random variable $Z$.
\end{theorem}

\begin{theorem}[Tendsto majority arbitrary coalition zero of subsqrt]
\lean{FABL.tendsto_majority_arbitraryCoalition_zero_of_subsqrt}
\statementsource{Analysis-of-Boolean-Functions-by-Ryan-ODonnell}{pdf-ffbe9c7b3306-p275-b002}
\label{FABL.tendsto_majority_arbitraryCoalition_zero_of_subsqrt}
\leanok
\uses{FABL.coalitionalInfluence, FABL.coalitionalInfluence_majority_eq_of_card, FABL.majority, FABL.majorityHeadCoalition, FABL.tendsto_majority_head_total_coalitionalInfluence_zero_of_subsqrt}
\difficulty{3}
% dataset-pipeline: stub -- prose pending build_informal_statements.py
Statement pending regeneration.
\end{theorem}

\begin{theorem}[Coalitional influence when coalition size dominates]
\lean{FABL.tendsto_majority_arbitraryCoalition_one_of_omega_sqrt}
\proofsource{Analysis-of-Boolean-Functions-by-Ryan-ODonnell}{pdf-ffbe9c7b3306-p275-b002}
\label{FABL.tendsto_majority_arbitraryCoalition_one_of_omega_sqrt}
\leanok
\uses{FABL.coalitionalInfluence, FABL.coalitionalInfluence_majority_eq_of_card, FABL.majority, FABL.majorityHeadCoalition, FABL.tendsto_majority_head_total_coalitionalInfluence_one_of_omega_sqrt}
\difficulty{3}
Let $k,m:\mathbb{N}\to\mathbb{N}$, and for each $q$ let $J_q$ be a set of $k(q)$
coordinates among $k(q)+m(q)$ coordinates. If $m(q)\to\infty$ and
$k(q)/\sqrt{k(q)+m(q)}\to\infty$, then the probability that this coalition can make
majority attain both possible outputs, for a uniformly random input, converges to $1$.
\end{theorem}
